\section{Contact geometry and further directions}
\label{sec:discussion}

The atlas connects three levels of description. A parameter determines a compact
attractor; its coding map records which symbolic points are identified; and its
position within a prescribed contact family determines which of those
identifications are forced. The exact examples show why these levels are useful
together. The quarter-turn attractor has four different semilinear presentations
with the same first-level pieces, while their agreement is established by
explicit recodings. At the critical rectangle its first-level intersection is a
segment, with different address multiplicities along that segment. On a whole marked binary zipper
surface, by contrast, the stable set consists precisely of the embedded
canonical curves. A single geometric contact label cannot describe all of these
situations.

The triangle family makes the universal radial bounds concrete. Its elementary
altitude dissection belongs to the classical geometry of self-similar tiles;
related two-piece planar tiles are studied in~\cite{NgaiSirventVeermanWang2000}.
Here the explicit normalization places such a family on the connectedness
boundary at every contraction weight, and the separated perturbation proves
sharpness in the full opposite--opposite chart. The opposite diagonal gives a
second exact model: its circle and real antennae produce a discontinuous radial
ceiling. Thus a smooth displayed sheet can be informative while missing genuine
features of the parameter boundary.

Several precise questions follow from these constructions. Near a critical
triangle or rectangle, which directions in the full coefficient space preserve
contact, produce separation, or create positive-area overlap? The formulas for
the exact families provide base parameters at which to study this local change.
For a fixed eventually periodic pair, the contact equation is algebraic; one
can ask when its regular branches organize nearby fixed-address fibres and when
additional contact pairs change their local structure. These questions concern
both the geometry of a fibre and the way fibres intersect.

The stable-set construction supplies a common language for the complex-tree
families developed since~\cite{Espigule2013,Espigule2018,Espigule2019a} and for
marked zippers. Calegari's island theorem gives a disconnected embedded locus
on one zipper surface~\cite{Calegari2026}. It is natural to seek other prescribed
contact families with disconnected stable sets, and to compare their forced
quotients and excess-contact fibres. Passing from a slice to a larger family
changes the forced relation, so any persistence statement should specify the
family and coding relation under consideration. The exterior-coordinate
comparison makes this distinction visible: the four selected stable slices
are subsets of their marked zipper surfaces in the four-dimensional atlas,
and colour retains their contraction weight after projection to three
dimensions. Intersections in the projection can therefore arise from
different weights; they do not themselves identify systems or contacts.

The companion atlas makes these questions accessible as reproducible
experiments. Finite radial surveys suggest regions to investigate; exact
contact relations provide positive examples; complete prefix covers and the
contact-gap estimate provide open sets of certified separation. The next
computational step is to adaptively refine those records around specified
contact fibres and their intersections. This extends the atlas through explicit
mathematical questions and checkable parameter data, while the visualization
continues to support immediate movement between parameters and attractors.
